\section{Methods}
\label{sec:methods}

Our analysis comprises the following steps:
\nopagebreak

\paragraph{Data Integration}

Substance records were compiled from 18 European regulatory and monitoring sources and one Flemish administrative list. These include ECHA obligation lists (Candidate List~\cite{ECHA_CandidateList}, Authorisation List~\cite{ECHA_AuthorisationList}, Restriction List~\cite{ECHA_RestrictionList}, POPs List~\cite{ECHA_POPsList}, EU Positive List~\cite{ECHA_EUPositiveList}, and the CLH Harmonised Classification and Labelling List~\cite{ECHA_CLHList}), ECHA activity lists (Restriction Process~\cite{ECHA_RestrictionProcess}, SVHC Identification~\cite{ECHA_SVHCIdentification}, Authorisation Process~\cite{ECHA_AuthorisationProcess}, Dossier Evaluation~\cite{ECHA_DossierEvaluation}, CLH Process~\cite{ECHA_CLHProcess}, Substance Evaluation~\cite{ECHA_SubstanceEvaluation}, POPs Process~\cite{ECHA_POPsProcess}, PBT Assessment~\cite{ECHA_PBTAssessment}, and ED Assessment~\cite{ECHA_EDAssessment}), as well as REACH registration data~\cite{ECHA_REACHRegistrations}. 
Additional sources comprise the EU Pesticides Database (Active Substances)~\cite{EU_Pesticides}, the OSPAR List of Chemicals for Priority Action~\cite{OSPAR_PriorityChemicals}, and a Flemish administrative substance list developed in the context of water quality modelling~\cite{Vlaanderen_SommatieStoffen}. This Flemish \emph{summation substances} list reflects the heterogeneity of \emph{substance names} in Flemish environmental permits and monitoring. 

For each entry, available identifiers (CAS number, EC number, and substance name) were collected. Where possible, substances were linked to structure-based identifiers (InChIKey) via PubChem~\cite{Kim2015}. The resulting dataset establishes a unified mapping between regulatory entries, their identifiers, and their originating source lists.

\paragraph{Chemical Identity and Linkability}
\label{sec:linkability}

Each substance entry was classified into one of several types reflecting its degree of chemical specificity: structure-defined molecule, substance group, mixture/UVCB, analytical parameter, or administrative entry.

Based on this classification, a seven-tier linkability taxonomy was defined, ranging from tier~1 (directly structure-resolvable via InChIKey) through partially resolvable entries (tiers~2--3), unresolved CAS numbers (tier~4), mixtures and UVCBs (tier~5), and administrative entries (tier~6), to tier~7 (no usable identifier). This taxonomy expresses the degree to which regulatory substances can be unambiguously linked across datasets.

\paragraph{Identifier Consistency Analysis}

For entries with both CAS numbers and InChIKeys, bidirectional mappings were analysed to assess identifier consistency. This analysis evaluates whether CAS numbers uniquely map to structures and vice versa.

\paragraph{Structure Coverage and Cross-List Overlap}

For each regulatory source, the proportion of structure-defined substances was calculated. Overlap between sources was analysed using structure-based matching (i.e.,~using InChIKey) and set similarity measures.

\paragraph{Analysis of Non-Structure Substances}

Substances without structure identifiers were analysed using text-based methods, including embedding-based clustering~\cite{Rousseeuw1987,Lloyd1982,McInnes2018} and similarity matching to chemical taxonomies (i.e.,~using ChemOnt), to assess whether semantic techniques can recover meaningful structure or classification~\cite{Manning2008,DjoumbouFeunang2016}. Each substance name and ChemOnt class label was encoded using the \path{all-MiniLM-L6-v2} sentence-transformer model~\cite{Reimers2019,Wang2020MiniLM,devlin-etal-2019-bert}, producing a 384-dimensional dense embedding vector per substance entry.

\paragraph{Embedding Model Validation Using ChemOnt Ground Truth}
\label{sec:chemont_validation}

To evaluate the ability of text-based embeddings to recover chemical classifications from substance names, an internal validation was performed using structure-defined substances as ground truth. Substances with an
unambiguous InChIKey were linked to the ChemOnt taxonomy through ClassyFire classification~\cite{DjoumbouFeunang2016}, providing reference assignments at the levels of direct parent, subclass, class, superclass, and kingdom.

Substance names and ChemOnt class labels were encoded using pre-trained transformer models (\path{all-MiniLM-L6-v2}, \path{all-mpnet-base-v2}~\cite{Song2020MPNet}, \path{allenai/scibert_scivocab_uncased}~\cite{Beltagy2019SciBERT}, and \path{dmis-lab/biobert-v1.1}~\cite{Lee2020BioBERT}) within the \emph{sentence-transformers} framework~\cite{Reimers2019,Wang2020MiniLM,devlin-etal-2019-bert}. Cosine similarity was used to rank all ChemOnt classes for each substance name.

Performance was assessed using Hit@k ($k = 1, 3, 5, 10$)~\cite{Manning2008}. To account for the hierarchical structure of ChemOnt, Hit@1 was additionally computed for each taxonomy level.

This validation provides a comparison of embedding models and demonstrates the extent to which semantic representations of substance names approximate structure-based chemical taxonomy assignments.

To assess whether large language models provide a more capable baseline, Claude Sonnet~4.6 was applied to the same validation benchmark under a zero-shot prompting protocol. For each substance name, the model was prompted to generate the ChemOnt class label at each hierarchy level (kingdom, superclass, class, subclass, and direct parent) together with a confidence estimate (high, medium, or low). Because the LLM generates a single predicted label per hierarchy level rather than ranking the full class space, only Hit@1 (accuracy) is reported; Hit@k for $k > 1$ is not defined for this setting. Performance metrics are therefore not directly comparable to the embedding retrieval results but address the same evaluation question.

\paragraph{LLM-Based Equivalence Assessment for Non-Structure Substances}

To identify regulatory group entries whose scope can be aligned with a ChemOnt taxonomy node, Claude Sonnet~4.6 was applied to a sample of non-structure substance names. For each entry, the model assessed whether the name describes a ChemOnt-equivalent group (rather than a single substance), identified the most specific equivalent ChemOnt node, and classified the relationship as \texttt{skos:exactMatch}, \texttt{skos:broadMatch} (regulatory entry is a subset of the ChemOnt class), or \texttt{skos:narrowMatch}. High-confidence assignments were serialised as SKOS triples and added to the knowledge graph.

\paragraph{Knowledge Graph Construction}

All data were integrated into an RDF knowledge graph using SKOS~\cite{W3C2009SKOS} for taxonomy representation. Structure-defined substances were linked to ChemOnt classes and to ChEBI entities via InChIKey, enabling integration of biological hazard annotations~\cite{W3C2017OA}.

\paragraph{Workload Estimation}

For non-structure substances, interoperability requires manual assessment of pairwise relationships (equivalence, subset, overlap, disjoint). The required effort scales quadratically with the number of such entries as $\binom{N}{2} = N(N-1)/2$, where $N$ is the number of non-structure entries under consideration.

\paragraph{Priority Scoring}
\label{sec:priority_score}

To illustrate the analytical value of the integrated knowledge graph (C5), a composite priority score was computed by combining regulatory list membership with biological role annotations from ChEBI~\cite{Degtyarenko2008}, following the ToxPi framework~\cite{Reif2010}. This score is illustrative only and does not constitute a formal risk assessment.